\documentclass[11pt]{article}
\usepackage[T1]{fontenc}
\usepackage{lmodern}
\usepackage[a4paper,margin=27mm]{geometry}
\usepackage{microtype,amsmath,amssymb,amsthm,mathtools}
\usepackage{booktabs,array,longtable,enumitem,xcolor}
\usepackage{listings}
\usepackage[colorlinks=true,linkcolor=blue!45!black,urlcolor=blue!55!black,citecolor=blue!45!black]{hyperref}
\usepackage{fancyhdr}
\hypersetup{pdftitle={Stirling formula in intuitionistic primitive recursive arithmetic},pdfauthor={Prepared for Vladimir Reshetnikov},pdfsubject={Finite rational proof, explicit error bounds, and Lean formalization plan}}
\pagestyle{fancy}\fancyhf{}
\fancyhead[L]{\small Stirling's formula in iPRA}
\fancyhead[R]{\small A finite rational proof}
\fancyfoot[C]{\thepage}
\setlength{\headheight}{14pt}
\setlength{\emergencystretch}{3em}
\setlist{itemsep=3pt,topsep=5pt}
\lstset{basicstyle=\ttfamily\footnotesize,breaklines=true,columns=fullflexible,keepspaces=true,frame=single,rulecolor=\color{black!20}}
\newtheorem{theorem}{Theorem}[section]
\newtheorem{lemma}[theorem]{Lemma}
\newtheorem{proposition}[theorem]{Proposition}
\newtheorem{corollary}[theorem]{Corollary}
\theoremstyle{definition}
\newtheorem{definition}[theorem]{Definition}
\newtheorem{remark}[theorem]{Remark}
\newcommand{\N}{\mathbb N}
\newcommand{\Q}{\mathbb Q}
\newcommand{\ipra}{\mathsf{iPRA}}
\newcommand{\eps}[1]{2^{-{#1}}}
\newcommand{\abs}[1]{\left|#1\right|}
\newcommand{\Log}{\operatorname{Log}}
\newcommand{\At}{\operatorname{At}}
\newcommand{\bit}{\operatorname{bit}}
\newcommand{\Rat}{\operatorname{Rat}}
\newcommand{\ArithCheck}{\operatorname{Check}}
\newcommand{\PR}{\mathsf{PR}}
\newcommand{\lmargin}{\lambda}
\newcommand{\umargin}{\mu}
\newcommand{\repo}{\texttt{329161f8da0b319d51ba9f6246c4b19d243f7193}}
\title{\textbf{Stirling's formula in intuitionistic\protect\\primitive recursive arithmetic}\\[6pt]
\large A complete finite rational argument, explicit error moduli,\\and a concrete path to Lean formalization}
\author{Prepared for Vladimir Reshetnikov}
\date{November 14, 2026}
\begin{document}
\maketitle
\begin{abstract}
We give a finitary proof of Stirling's formula in intuitionistic primitive
recursive arithmetic. Its numerical conclusion is stronger than the usual
asymptotic statement: for every integer $n\ge1$,
\[
 n!=\sqrt{2\pi n}\,(n/e)^n e^{r_n},
 \qquad \frac1{12n+1}<r_n<\frac1{12n}.
\]
The literal arithmetic theorem is an explicitly specified rational
inequality with a universally quantified precision parameter. Logarithms,
$\pi$, $e$, square roots, and the displayed real notation are interpreted by
fixed primitive recursive rational approximation algorithms. A rational
parametrization of the circle replaces trigonometric integrals, and finite
telescoping estimates replace integration by parts and convergence
principles. In particular, no completeness, choice, unbounded search, or
induction on a formula with an unbounded quantifier is used.

The written argument and its arithmetic reduction are supplied in full.
The accompanying Lean file is explicitly uncompiled scaffolding: it gives
concrete rational programs, the exact object-language target in ProveIt's
\texttt{ConciseIPRA}, and the transfer to its native calculus, but does not
claim a completed machine-checked derivation. Executable exact-rational
checks accompany the proof and serve only as supplementary validation.
\end{abstract}
\newpage
\tableofcontents
\newpage

\section{The result and its precise foundational meaning}
\label{sec:statement}

The familiar statement is
\begin{equation}\label{eq:stirling}
 \frac{n!}{\sqrt{2\pi n}(n/e)^n}\longrightarrow1.
\end{equation}
An arithmetic theory has natural-number variables, not a sort of arbitrary
real numbers. Consequently, saying that \eqref{eq:stirling} is provable in
$\ipra$ requires both a representation and an effective assertion about it.
Here every analytic expression names a particular rational algorithm with
proved error bounds. Nothing asserts that all real numbers or all
convergent sequences exist in the theory.

\begin{theorem}[Quantitative Stirling formula]\label{thm:main}
In the interpretation by explicit primitive recursive rational names
developed below, $\ipra$ proves the following for every $n\ge1$.
Put
\begin{align}
 a_n&=\sum_{j=1}^n\log j-(n+\tfrac12)\log n+n,\label{eq:a}\\
 C&=\tfrac12\log(2\pi),\qquad r_n=a_n-C,\label{eq:r}\\
 R_n&=\frac{n!e^n}{n^n\sqrt{2\pi n}}.\label{eq:R}
\end{align}
Then
\begin{equation}\label{eq:robbins}
 \frac1{12n+1}<r_n<\frac1{12n},\qquad \log R_n=r_n,
\end{equation}
and, in particular,
\begin{equation}\label{eq:relative}
 0<R_n-1<\frac1{12n-1}.
\end{equation}
The function $N(p)=2^p$ is a primitive recursive modulus for
\eqref{eq:stirling}: $n\ge2^p$ implies $\abs{R_n-1}<2^{-p}$.
\end{theorem}

The constants in \eqref{eq:robbins} are the classical Robbins bounds
\cite{Robbins}. The purpose here is their explicit realization in the
restricted intuitionistic arithmetic, including a finite proof of the
normalizing constant. There is no claim of a new numerical bound.

For clarity, here is the exact arithmetic endpoint. Sections
\ref{sec:log} and \ref{sec:pi} define rational primitive recursive functions
$L(q,t)$ and $P(t)$ approximating $\log q$ and $\pi$. For $n\ge1$ set
\begin{align}
 b(n)&=\min\{b\le n+1:n<2^b\},& t(n,p)&=p+b(n)+4,\label{eq:precision}\\
 Q(n,p)&=\sum_{j=1}^nL(j,t(n,p))
 -(n+\tfrac12)L(n,t(n,p))+n\notag\\
 &\hspace{25mm}-\tfrac12 L(2P(t(n,p)),t(n,p)).\label{eq:Q}
\end{align}
Put $Q(0,p)=0$ to make the function total. A rational is represented by
signed integer numerator and positive integer denominator, so the following
is an ordinary first-order arithmetic sentence, with a quantifier-free
matrix after the indicated primitive recursive definitions are expanded:
\begin{equation}\label{eq:internal}
 \boxed{\quad
 \forall n,p\quad n\ge1\ \Longrightarrow\
 \frac1{12n+1}-2^{-p}\le Q(n,p)
 \le\frac1{12n}+2^{-p}.\quad}
\end{equation}
The comparisons mean cross-multiplied integer comparisons, not real-order
predicates. We prove this sentence and, more strongly, its version with
the positive margins in \eqref{eq:margins} below. This is the literal
$\ipra$ theorem behind \eqref{eq:robbins}.

\paragraph{What is finished.}
This manuscript provides a human-readable proof, explicit definitions,
finite estimates for every analytic step, and the induction audit. The
archive also contains executable rational algorithms and a Lean plan tied
to inspected repository APIs. It does \emph{not} contain a completed Lean
term inhabiting the new object-language derivability judgment. A precise
description of that remaining implementation appears in
Section~\ref{sec:lean}.

\section{Arithmetic conventions and finite notation}
\label{sec:arith}

We use the usual first-order presentation of $\ipra$: intuitionistic
predicate logic with equality, symbols for primitive recursive functions,
their defining equations, successor separation and injectivity, and
induction for quantifier-free formulas. This is the presentation in
ProveIt's \texttt{ConciseIPRA/PR.lean} \cite{ProveItPR}.

Integers can be pairs of naturals, with $(u,v)$ representing $u-v$.
Rationals can be triples $(u,v,d)$ representing $(u-v)/d$, with $d>0$.
Addition and multiplication use cross products, and order uses integer
cross multiplication. Lowest-term normalization is unnecessary. Every
function is made total by assigning a fixed default on malformed inputs;
the proofs maintain validity of the inputs actually used.

Finite sums, products, exponentiation, factorials, bounded maxima, and
bounded searches are primitive recursive. All rational comparisons used in
searches are decidable integer comparisons. Finite lists of partition
points can be encoded by primitive recursive pairing or list encodings;
insertion sort uses a loop bounded by the list length. These are finite
algorithms with quantifier-free correctness invariants.

\paragraph{Named quantities.}
The notation $x=\log q$, for example, abbreviates a specified sequence of
rational enclosures. Equality of two such names means that any prescribed
rational error tolerance is met by explicit finite approximations. We
never introduce a new existentially chosen limit. To combine a finite
number of estimates, choose all precisions at the start and add their
error budgets. For a bounded sum with $s$ summands, a summand error
$\delta/(s+1)$ gives total error at most $\delta$.

This notation must not be confused with permission to induct on
``the equality holds at every precision.'' In every inductive argument
below, the precision and final finite bounds are fixed parameters. The
induction predicate is a rational identity or inequality at those fixed
parameters. A bounded conjunction can be encoded by a primitive recursive
Boolean product. Section~\ref{sec:elimination} gives a more explicit audit.

\section{Finite rational sums in place of integration}
\label{sec:sums}

Only a small amount of elementary ``calculus'' is needed. We derive its
finite content now. An integral symbol below is an abbreviation for these
sums and their error estimates.

\subsection{Refinement, concatenation, and a modulus}

For a rational partition $a=t_0<\cdots<t_N=b$ with $0<a<b$, define
\[
 U(\mathcal P)=\sum_{i<N}\frac{t_{i+1}-t_i}{t_i},\qquad
 D(\mathcal P)=\sum_{i<N}\frac{t_{i+1}-t_i}{t_{i+1}}.
\]
Splitting a cell at $w\in[u,v]$ increases $D$ and decreases $U$:
this follows by multiplying through by the positive denominators and
using $u\le w\le v$. Induction on the number of added points proves the
same for any refinement. Two finite rational partitions have their sorted
union as a common refinement. Their intervals $[D,U]$ therefore overlap.

For the uniform partition with $h=(b-a)/N$, the width is exactly
\begin{equation}\label{eq:logwidth}
 U-D=h\left(\frac1a-\frac1b\right).
\end{equation}
Thus, for a positive rational $x=A/B$, with $H=A+B$, the uniform
partition of $[\min(1,x),\max(1,x)]$ into $N=2^rH^4$ pieces has
width at most $2^{-r}$: use $b-a\le H$, $a\ge1/H$, and
$U-D\le(b-a)^2/(Na^2)$.

Oriented partitions handle a reversed interval by changing its sign.
Concatenating two partitions adds their sums. When the middle endpoint
lies outside the other two, the overlap cancels after a common refinement.
There are only finitely many rational endpoint orderings, each decidable.
No assertion about an infinite collection of partitions is required.

More generally, for any of the explicitly given rational functions $f$
used here, suppose a proved rational bound is
$\abs{f(u)-f(v)}\le B\abs{u-v}$ on a rational interval of length $d$.
The left sum $S_N$ on a uniform partition and its refinement $S_{NM}$
satisfy
\begin{equation}\label{eq:refine}
 \abs{S_N-S_{NM}}\le\frac{Bd^2}{N},\qquad
 \abs{S_N-S_M}\le Bd^2\left(\frac1N+\frac1M\right).
\end{equation}
Indeed, each sample point moves by at most $d/N$ inside its original cell;
multiply this variation bound by the cell length and sum. This is a finite
double-sum proof. It supplies a primitive recursive modulus whenever we
write an integral of one of these functions.

\subsection{Polynomial sums and finite telescoping}

For $m\ge0$ and $0\le u\le v$, factorization gives
\[
 v^{m+1}-u^{m+1}=(v-u)\sum_{j=0}^m v^{m-j}u^j.
\]
Each summand lies between $u^m$ and $v^m$. Apply this on every partition
cell and telescope. The lower and upper sums for $t^m$ enclose
$(b^{m+1}-a^{m+1})/(m+1)$. On $[0,x]$ with a uniform partition their
difference is at most $x^{m+1}/N$. In the sum notation this proves
\begin{equation}\label{eq:monomial}
 \int_0^x t^m\,dt=\frac{x^{m+1}}{m+1},
\end{equation}
including its explicit rational error. Finite linear combinations give the
same statement for every particular polynomial used in the argument.

Another frequently used finite principle is the following. If rational
functions $G,F$ satisfy
\[
 \abs{G(v)-G(u)-(v-u)F(u)}\le B(v-u)^2
 \quad(a\le u\le v\le b),
\]
then summation over a uniform mesh yields
\begin{equation}\label{eq:FTCfinite}
 \abs{\frac{b-a}{N}\sum_{j<N}F(a+j(b-a)/N)
       -(G(b)-G(a))}\le\frac{B(b-a)^2}{N}.
\end{equation}
This is merely the triangle inequality and telescoping; a fundamental
theorem of calculus is not assumed. We will exhibit the rational
increment identities needed to apply it.

\section{Logarithms as explicit rational algorithms}
\label{sec:log}

\subsection{The logarithm laws}

For rational $x>0$, let $\log x$ name the oriented enclosures from
Section~\ref{sec:sums} on the interval from $1$ to $x$. Scaling every
partition point by a positive rational $x$ preserves each summand:
\[
 \frac{xu_{i+1}-xu_i}{xu_i}
 =\frac{u_{i+1}-u_i}{u_i}.
\]
Concatenation at $x$ therefore proves
\begin{equation}\label{eq:loglaws}
 \log(xy)=\log x+\log y,\quad \log1=0,\quad
 \log(1/x)=-\log x.
\end{equation}
More explicitly, refine the three partitions involved to common scaled
partitions. The discrepancy of their original sums is at most the sum of
their three enclosure widths. This proves every prescribed finite-
precision instance of \eqref{eq:loglaws}.

Bounding $1/t$ by its endpoint values gives, for $0<x<y$,
\begin{equation}\label{eq:loglip}
 \frac{y-x}{y}\le\log y-\log x\le\frac{y-x}{x}.
\end{equation}
In particular, with the orientation treated correctly on either side of
$1$, for all $x>0$ one has
\begin{equation}\label{eq:logglobal}
 1-\frac1x\le\log x\le x-1.
\end{equation}
For rational endpoints these are finite-sum comparisons. For a particular
named positive argument bounded away from zero, the same inequalities
follow with its explicit input error added. The Lipschitz estimate in
\eqref{eq:loglip} controls that addition. No case split on an undecidable
comparison of real numbers is needed.

\subsection{A finite positive series and its tail}

For rational $0\le z<1$, translation of partitions and
\eqref{eq:loglaws} give
\[
 \log\frac{1+z}{1-z}
 =\int_0^z\left(\frac1{1+t}+\frac1{1-t}\right)dt
 =2\int_0^z\frac{dt}{1-t^2}.
\]
For each finite $K$, the exact geometric identity is
\[
 \frac1{1-t^2}=\sum_{j<K}t^{2j}+\frac{t^{2K}}{1-t^2}.
\]
On $[0,z]$ the last term is nonnegative and at most
$t^{2K}/(1-z^2)$. Multiply by partition widths, sum, and use
\eqref{eq:monomial}. We obtain the explicit enclosure
\begin{equation}\label{eq:atanh}
 0\le\log\frac{1+z}{1-z}-2\sum_{j<K}\frac{z^{2j+1}}{2j+1}
 \le\frac{2z^{2K+1}}{(2K+1)(1-z^2)}.
\end{equation}
This is a finite geometric calculation with a tail bound, not an appeal
to an infinite-series theorem.

\subsection{The algorithm \texorpdfstring{$L(q,r)$}{L(q,r)} used in the arithmetic statement}

For $q=A/B\ge1$, find by bounded search the largest $k\le A$ with
$2^kB\le A$. Then $2^k\le q<2^{k+1}$. Define
\[
 y=q/2^k,\quad z=(y-1)/(y+1),\quad K=r+k+4,
\]
so $1\le y<2$ and $0\le z<1/3$. Put
\begin{equation}\label{eq:L}
 L(q,r)=2k\sum_{j<K}\frac{3^{-(2j+1)}}{2j+1}
       +2\sum_{j<K}\frac{z^{2j+1}}{2j+1}.
\end{equation}
For $0<q<1$, set $L(q,r)=-L(1/q,r)$, and set $L(q,r)=0$ on
invalid nonpositive input. The latter case will never be used in a valid
logarithmic calculation.

By \eqref{eq:atanh}, each of the two unweighted tails in
\eqref{eq:L} is at most $(3/4)9^{-K}$. Therefore
\begin{equation}\label{eq:Lerr}
 \abs{L(q,r)-\log q}
 \le\tfrac34(k+1)9^{-(r+k+4)}\le2^{-r}.
\end{equation}
For example, the last inequality follows from $k+1\le9^k$ and
$2^r\le9^r$, both proved by quantifier-free induction. There is large
slack in this estimate. The actual geometric tail can be retained until
the last comparison with a rational enclosure, so no exact rational
inequality is obtained by an illicit limit argument.

The search bound is legitimate: $2^A>A\ge q$ for positive integer $A$
unless the result already occurs earlier, so the maximal exponent is
found within the declared finite range. All operations in \eqref{eq:L}
are primitive recursive.

\section{An independent primitive recursive name for \texorpdfstring{$\pi$}{pi}}
\label{sec:pi}

First define $\At(x)$, for rational $0\le x\le1$, by the finite sums
\[
 A_N(x)=\frac{x}{N}\sum_{j<N}\frac1{1+(xj/N)^2}.
\]
The summand decreases. The corresponding left-right gap is
$\frac{x}{N}(1-1/(1+x^2))\le1/N$. Thus these are explicitly
compatible rational enclosures. Define $\pi=4\At(1)$. This definition
does not use factorials or Stirling's formula; a direct geometric
identification is given below.

The finite identity
\[
 \frac1{1+t^2}=\sum_{j<K}(-1)^jt^{2j}
                +\frac{(-1)^Kt^{2K}}{1+t^2}
\]
and \eqref{eq:monomial} imply
\begin{equation}\label{eq:atan}
 S_K(x)=\sum_{j<K}\frac{(-1)^jx^{2j+1}}{2j+1},\qquad
 \abs{\At(x)-S_K(x)}\le\frac{x^{2K+1}}{2K+1}.
\end{equation}

\subsection{A fast formula, proved by finite substitution}

We need only the particular addition law
$\At(1)=\At(1/2)+\At(1/3)$. Here is its finite proof. Put
$a=1/2$, $b=1/3$, $h=b/N$ and
\[
 \phi(u)=\frac{a+u}{1-au},\qquad t_j=\phi(jh),\qquad
 f(t)=\frac1{1+t^2}.
\]
Then $\phi(0)=a$, $\phi(b)=1$, and direct rational algebra gives
\[
 \phi(v)-\phi(u)=\frac{(1+a^2)(v-u)}{(1-au)(1-av)},\qquad
 f(\phi(u))=\frac{(1-au)^2}{(1+a^2)(1+u^2)}.
\]
Consequently the left sum $T_N$ on the partition $t_j$ of $[a,1]$
satisfies
\[
 (t_{j+1}-t_j)f(t_j)
 =h f(jh)\frac{1-ajh}{1-a(j+1)h}.
\]
Its difference from the corresponding summand of $A_N(b)$ is between
$0$ and $(3/5)h^2$. Thus $0\le T_N-A_N(b)\le1/(15N)$.
Also $\max_j(t_{j+1}-t_j)\le3/(5N)$, so the left-right gap of
$T_N$ is at most $9/(50N)$. The quantity
$A_{2N}(1)-A_N(a)$ is exactly the uniform left sum on $[a,1]$ and
has gap $3/(20N)$. Compare the partitions by their finite common
refinement. Adding the three displayed errors gives
\begin{equation}\label{eq:atanaddfinite}
 \abs{A_{2N}(1)-A_N(1/2)-A_N(1/3)}\le\frac1N.
\end{equation}
This establishes the named equality with an explicit modulus.

We use the following efficient rational code for $\pi$:
\begin{equation}\label{eq:Pi}
 P(r)=4\bigl(S_{r+4}(1/2)+S_{r+4}(1/3)\bigr).
\end{equation}
If $K=r+4$, its error is at most
\begin{equation}\label{eq:Pierr}
 \frac{4(2^{-2K-1}+3^{-2K-1})}{2K+1}
 \le2^{2-2K}\le2^{-r-6}<2^{-r}.
\end{equation}
Alternating finite sums, grouped into successive pairs, give for $K\ge2$
\[
 \frac{505}{162}
 =4\left(\frac12-\frac1{24}+\frac13-\frac1{81}\right)
 \le4(S_K(1/2)+S_K(1/3))\le\frac{10}{3}.
\]
Hence $3<P(r)<4$ and $3<\pi<4$. These explicit bounds will keep
denominators and logarithmic arguments safely positive.

\subsection{Why this is the ordinary circle constant}
\label{sec:circle}

For $0\le t\le1$, define rational functions
\begin{equation}\label{eq:scw}
 s(t)=\frac{2t}{1+t^2},\quad
 c(t)=\frac{1-t^2}{1+t^2},\quad
 w(t)=\frac2{1+t^2}.
\end{equation}
Then $c(t)^2+s(t)^2=1$, and $X(t)=(c(t),s(t))$ parametrizes the
first quadrant of the unit circle; its inverse is $t=y/(1+x)$.
The exact chord identity is
\[
 \abs{X(v)-X(u)}^2
 =\frac{4(v-u)^2}{(1+u^2)(1+v^2)}.
\]
Since $1+u^2\le1+v^2$, its nonnegative square root lies between
$(v-u)w(v)$ and $(v-u)w(u)$. Therefore the polygonal quarter-circle
length on the grid $j/N$ lies between
\[
 J_0(N)-\frac1N\quad\hbox{and}\quad J_0(N),\qquad
 J_0(N)=\frac1N\sum_{j<N}w(j/N)=2A_N(1).
\]
These bounds also apply after common refinements, giving the usual arc
length with an explicit modulus. The circumference divided by the
diameter is therefore $2\lim J_0(N)=4\At(1)=\pi$.
Square roots in the chord lengths can themselves be approximated by
bounded dyadic searches as in Section~\ref{sec:exp}. This paragraph
identifies the familiar constant; the arithmetic proof itself needs only
the rational sums and brackets just displayed.

\section{Wallis' inequality with a completely finite error term}
\label{sec:wallis}

The classical use of Wallis' product to determine the Stirling constant is
standard \cite{Bruin}. We need a version whose proof is a rational
calculation. The functions in \eqref{eq:scw} provide it.

Put $F_k(t)=w(t)s(t)^k$ and, for $N\ge1$,
\begin{equation}\label{eq:J}
 J_k(N)=\frac1N\sum_{j<N}F_k(j/N).
\end{equation}
We have $0\le s,c\le1$, $1\le w\le2$, and $c^2+s^2=1$.
For $0\le u\le v\le1$, with $h=v-u$ and $D(t)=1+t^2$,
\begin{align}
 s(v)-s(u)&=\frac{2h(1-uv)}{D(u)D(v)},\label{eq:sinc}\\
 c(v)-c(u)=w(v)-w(u)&=\frac{-2h(u+v)}{D(u)D(v)}.\label{eq:cinc}
\end{align}
In particular $s$ increases, $c,w$ decrease, $\abs{s(v)-s(u)}\le2h$,
and $\abs{c(v)-c(u)}\le4h$. Factoring powers gives
$\abs{b^k-a^k}\le k\abs{b-a}$ for $a,b\in[0,1]$, and consequently
\[
 \abs{F_k(v)-F_k(u)}\le4(k+1)h.
\]
Thus \eqref{eq:refine} supplies explicit moduli for every $J_k$.
We may write $I_k$ for these names; in particular $I_0=\pi/2$.

\subsection{The finite recurrence}

We will prove, entirely over the rationals,
\begin{equation}\label{eq:Jrec}
 \abs{(k-1)J_{k-2}(N)-kJ_k(N)}\le\frac{8k^2}{N}\quad(k\ge2),
 \qquad \abs{J_1(N)-1}\le\frac{12}{N}.
\end{equation}

Subtract the proposed linear increments from
\eqref{eq:sinc}--\eqref{eq:cinc}. Direct denominator clearing gives
\begin{align*}
 s(v)-s(u)-hc(u)w(u)
 &=\frac{-2h^2[u(3-u^2)+(1-u^2)h]}{D(u)^2D(v)},\\
 c(v)-c(u)+hs(u)w(u)
 &=\frac{h^2[-2+6u^2+4uh]}{D(u)^2D(v)}.
\end{align*}
Since all denominators are at least $1$ and $u,h\in[0,1]$, the
absolute values are at most $8h^2$ and $12h^2$, respectively.

For $0\le a\le b\le1$ and integer $m\ge1$, telescoping powers
twice gives
\begin{equation}\label{eq:powerr}
 0\le b^m-a^m-ma^{m-1}(b-a)
 \le\frac{m(m-1)}2(b-a)^2.
\end{equation}
To see the upper bound, factor $b^m-a^m$ once, subtract its $m$
copies of $a^{m-1}$, and apply
$b^j-a^j\le j(b-a)$ to each remaining term. Summing the integers
$0,1,\ldots,m-1$ proves \eqref{eq:powerr}; this also handles $m=1$.

Apply it to $a=s(u)$, $b=s(v)$. The error on replacing
$s(v)^m-s(u)^m$ by $hm s(u)^{m-1}c(u)w(u)$ is at most
$(2m^2+6m)h^2$. Set
\[
 G_k(t)=c(t)s(t)^{k-1},\qquad
 H_k(t)=(k-1)F_{k-2}(t)-kF_k(t).
\]
The product increment for $G_k$ has three error terms: the power error
$2m^2+6m$, the $c$ error $12$, and the cross term $8m$, all
multiplied by $h^2$, with $m=k-1$. Its linear term is $hH_k(u)$,
by $c^2+s^2=1$. Therefore
\[
 \abs{G_k(v)-G_k(u)-hH_k(u)}
 \le2k(k+5)h^2\le8k^2h^2\qquad(k\ge2).
\]
Now telescope on the grid $j/N$. Since $G_k(0)=G_k(1)=0$,
\eqref{eq:FTCfinite} proves the first assertion of
\eqref{eq:Jrec}. Telescoping the increment for $c$, using $c(0)=1$
and $c(1)=0$, proves the second. These calculations contain the full
finite replacement for integration by parts.

\subsection{The rational Wallis bounds}

Define rational products
\begin{align}
 p_m&=\prod_{j=1}^m\frac{2j-1}{2j}
     =\frac{\binom{2m}{m}}{4^m},&
 q_m&=\prod_{j=1}^m\frac{2j}{2j+1},\notag\\
 W_m&=\frac{q_m}{p_m}
     =\prod_{j=1}^m\frac{4j^2}{4j^2-1}
     =\frac1{(2m+1)p_m^2}.\label{eq:W}
\end{align}
The displayed equalities are factorial cancellations, proved by finite
product induction. Divide the recurrence error in \eqref{eq:Jrec} by
$k$ and iterate it. The multiplicative coefficients $(k-1)/k$ are at
most $1$, so the errors accumulate to give
\begin{align}
 \abs{J_{2m}(N)-p_mJ_0(N)}&\le\frac{8m(m+1)}N,\label{eq:evenerr}\\
 \abs{J_{2m+1}(N)-q_m}&\le\frac{8m^2+16m+12}N.\label{eq:odderr}
\end{align}
For example the first error is $\sum_{j=1}^m16j/N$; the second
starts with $12/N$ and adds $\sum_{j=1}^m8(2j+1)/N$.
These are quantifier-free inductions with $N$ retained as a parameter.

Termwise $0\le s\le1$ gives
\[
 J_{2m+1}(N)\le J_{2m}(N)\le J_{2m-1}(N).
\]
Also $p_m\ge2^{-m}$, because every factor is at least $1/2$.
Combining these inequalities with \eqref{eq:evenerr}--\eqref{eq:odderr}
proves the following wholly rational theorem.

\begin{theorem}[Finite Wallis inequality]\label{thm:finitewallis}
For every $m,N\ge1$, put $D_m=16\,2^m(m+1)^2$. Then
\begin{equation}\label{eq:finitewallis}
 \boxed{\quad W_m-\frac{D_m}N\le J_0(N)
 \le\frac{2m+1}{2m}W_m+\frac{D_m}N.\quad}
\end{equation}
\end{theorem}

\begin{proof}
For the lower bound, the sum of the two errors before division by $p_m$
is $(16m^2+24m+12)/N\le16(m+1)^2/N$. For the upper bound, use
the odd estimate with index $m-1$. The error sum is
$(16m^2+8m+4)/N\le16(m+1)^2/N$; its main term is
$q_{m-1}/p_m=((2m+1)/(2m))W_m$. Divide by $p_m\ge2^{-m}$.
\end{proof}

Since $J_0(N)$ names $\pi/2$ with a left-right gap $1/N$,
the named-code reading is
\begin{equation}\label{eq:wallis}
 W_m\le\frac\pi2\le\frac{2m+1}{2m}W_m.
\end{equation}
This notation hides no convergence principle: at tolerance $2^{-r}$,
choose, for example, $N=2^{r+2}(D_m+1)$ in
\eqref{eq:finitewallis}. All errors are then explicitly bounded.

One useful strictly rational bound is $W_m\le2$. Write $W_m=A/B$
with $B\ge1$, use $J_0(N)\le2$, and substitute $N=2D_mB$ in
\eqref{eq:finitewallis}. Then $2A\le4B+1$; discreteness of
integers gives $A\le2B$. This illustrates how an exact rational
conclusion is obtained by a finite substitution, without taking an
unjustified limit.

\section{The logarithmic factorial estimate}
\label{sec:factorial}

We now prove the numerical core of the theorem. In this section all named
equalities are the finite-precision identities established above. Their
elimination to \eqref{eq:internal} is made explicit in
Section~\ref{sec:elimination}.

\subsection{A positive difference with two rational telescopers}

For $n\ge1$ let $a_n$ be as in \eqref{eq:a}. Cancellation and the
logarithm addition law give
\begin{equation}\label{eq:d}
 d_n:=a_n-a_{n+1}
 =\left(n+\tfrac12\right)\log(1+1/n)-1.
\end{equation}
Put $z=1/(2n+1)$. Then $(1+z)/(1-z)=(n+1)/n$ and
$n+1/2=1/(2z)$. Dividing the finite formula \eqref{eq:atanh} by
$2z$ and subtracting its first term gives
\begin{equation}\label{eq:dseries}
 d_n=\sum_{j\ge1}\frac{z^{2j}}{2j+1}.
\end{equation}
The sum symbol denotes finite partial sums and their geometric tails.
For example, after $j=K$ the tail is at most
$z^{2K+2}/((2K+3)(1-z^2))$.

Its first term is $z^2/3$. Comparing every denominator with $3$ and
retaining the improvement at $j=2$ gives
\begin{equation}\label{eq:dcompare}
 \frac1{3(2n+1)^2}\le d_n
 \le\frac{z^2}{3(1-z^2)}-\frac{2z^4}{15}
 =\frac1{12n(n+1)}-\frac{2}{15(2n+1)^4}.
\end{equation}
For completeness, the coefficient gap at $j=2$ is
$1/3-1/5=2/15$; every later gap is nonnegative. This proves the
upper inequality at every finite precision.

Define
\begin{align}
 g_n&=\frac1{12n+1},& h_n&=\frac1{12n},\notag\\
 \lmargin_n&=\frac{24n-23}{3(2n+1)^2(12n+1)(12n+13)},&
 \umargin_n&=\frac{2}{15(2n+1)^4}.\label{eq:margins}
\end{align}
Both margins are positive for $n\ge1$. The exact polynomial identity
\[
 (12n+1)(12n+13)-36(2n+1)^2=24n-23
\]
shows that
\[
 \frac1{3(2n+1)^2}=g_n-g_{n+1}+\lmargin_n.
\]
Thus \eqref{eq:dcompare} becomes
\begin{equation}\label{eq:telescopers}
 g_n-g_{n+1}+\lmargin_n\le d_n
 \le h_n-h_{n+1}-\umargin_n.
\end{equation}

For any finite $m>n$, sum these inequalities. The logarithmic differences
telescope; discard all positive margins except the first:
\begin{equation}\label{eq:strongtel}
 g_n-g_m+\lmargin_n\le a_n-a_m
 \le h_n-h_m-\umargin_n.
\end{equation}
In particular, including $m=n$ separately,
\begin{equation}\label{eq:weaktel}
 0\le a_n-a_m\le\frac1{12n}-\frac1{12m}\qquad(m\ge n).
\end{equation}
This is finite telescoping, proved by quantifier-free induction after a
common precision is fixed. It does not assert convergence of an
arbitrary bounded monotone sequence.

\subsection{Determining the constant without choosing a limit}

Recall the independently defined $C=\frac12\log(2\pi)$.
The factorial identity in \eqref{eq:W}, together with logarithm addition,
gives
\begin{equation}\label{eq:F}
 F_m:=2a_m-a_{2m}
 =\frac12\log\left[(4+2/m)W_m\right].
\end{equation}
For a direct check, twice the expression on the left is
\[
 4\log(m!)-2\log((2m)!)+(4m+1)\log2-\log m,
\]
and $W_m=16^m(m!)^4/((2m+1)((2m)!)^2)$.

Wallis' bounds \eqref{eq:wallis} imply
\[
 2\pi\le(4+2/m)W_m\le2\pi\left(1+\frac1{2m}\right).
\]
Use the increasing logarithm and $\log(1+x)\le x$ to obtain
\begin{equation}\label{eq:Fbound}
 C\le F_m\le C+\frac1{4m}.
\end{equation}
On the other hand \eqref{eq:weaktel}, with endpoints $m,2m$, gives
$0\le a_m-a_{2m}\le1/(24m)$. Since
$a_m=F_m-(a_m-a_{2m})$, we have the explicit terminal estimate
\begin{equation}\label{eq:terminal}
 \boxed{\qquad -\frac1{24m}\le a_m-C\le\frac1{4m}.\qquad}
\end{equation}
No unknown constant or unspecified stage has been introduced.

Combining \eqref{eq:strongtel} with \eqref{eq:terminal}, for every
$m>n$, yields
\begin{equation}\label{eq:finitefinal}
 g_n+\lmargin_n-\frac1{8m}\le a_n-C
 \le h_n-\umargin_n+\frac1{6m}.
\end{equation}
For the lower error use $g_m\le1/(12m)$; the upper error is
$1/(4m)-1/(12m)=1/(6m)$. Given any requested tolerance, set
$m=\max(n+1,2^{p+8})$. The two residual errors are then smaller
than that tolerance by a known factor. In the rational-name notation this
proves
\begin{equation}\label{eq:strongrobbins}
 g_n+\lmargin_n\le r_n\le h_n-\umargin_n.
\end{equation}
Because the margins in \eqref{eq:margins} are explicit positive
rationals, this proves the strict Robbins inequalities
\eqref{eq:robbins} constructively. Strictness is not inferred merely
from a sequence of strict inequalities whose limit might be an equality.

\section{Returning to the usual factorial ratio}
\label{sec:exp}

It is possible to finish without developing a general exponential
function. We need only the ordinary constant $e$, square roots of a
specified positive name, and the global logarithm inequalities.

\subsection{A finite proof that \texorpdfstring{$\log e=1$}{log e = 1}}

For $m\ge1$ set
\[
 E_m=\sum_{j=0}^m\frac1{j!},\qquad
 B_m=(1+1/m)^m.
\]
Factorial comparison and a geometric sum give
\begin{equation}\label{eq:etail}
 0\le e-E_m\le\frac1{m\,m!}\le2^{1-m},
 \qquad 2\le E_m\le3.
\end{equation}
Here $e$ is the name supplied by this explicit tail estimate. To prove
the first bound, each successive term after $1/(m+1)!$ has ratio
at most $1/(m+1)$, so its geometric majorant sums to $1/(m\,m!)$.
The inequality $E_m\le3$ follows from $j!\ge2^{j-1}$ for $j\ge1$.
All these bounds hold already for finite tails.

The finite binomial theorem gives
\[
 B_m=\sum_{j=0}^m\frac1{j!}
             \prod_{i=0}^{j-1}(1-i/m).
\]
For $u_i\in[0,1]$, finite induction proves
$1-\sum_i u_i\le\prod_i(1-u_i)\le1$. Therefore
\begin{equation}\label{eq:EB}
 0\le E_m-B_m\le\frac1{2m}\sum_{j=2}^m\frac1{(j-2)!}
 \le\frac3{2m},\qquad 2\le B_m\le E_m\le3.
\end{equation}
The finite binomial theorem itself follows by multiplication and Pascal's
identity, using quantifier-free induction on its finite coefficient
array; no assertion about infinite power series occurs here.

Apply \eqref{eq:logglobal} to $1+1/m$ and multiply by $m$:
\[
 \frac{m}{m+1}\le\log B_m\le1.
\]
By \eqref{eq:loglip} and $B_m\ge2$,
$0\le\log E_m-\log B_m\le3/(4m)$. Hence
\begin{equation}\label{eq:loge}
 1-\frac1{m+1}\le\log E_m\le1+\frac3{4m}.
\end{equation}
Equations \eqref{eq:etail} and \eqref{eq:loge}, with a chosen
$m\ge2^{p+2}$ if tolerance $2^{-p}$ is wanted, prove
$\log e=1$ by explicit rational approximation. The series definition
identifies this with the usual Euler constant.

\subsection{Square roots and a direct rational ratio algorithm}

For a positive rational $a\ge1$ and integer $t\ge0$, define
\begin{equation}\label{eq:sqrtcode}
 T(a,t)=2^{-t}\left\lfloor\sqrt{\lfloor a2^{2t}\rfloor}\right\rfloor.
\end{equation}
The integer square root means the largest integer $j$ in a specified
finite range with $j^2\le\lfloor a2^{2t}\rfloor$. Searching up to
$\lfloor a2^{2t}\rfloor+1$ is a primitive recursive implementation.
Then the exact certificates are
\[
 T(a,t)^2\le a<(T(a,t)+2^{-t})^2.
\]
So \eqref{eq:sqrtcode} is a square-root approximation without a search
over real comparisons. Factoring $u^2-v^2=(u-v)(u+v)$ proves stability:
on inputs at least $1$, a change of at most $\delta$ in the input changes
the positive root by at most $\delta$.

For an explicit approximation of $R_n$ itself, set
\begin{align}
 t&=p+4n+2,& a&=2nP(t),& S&=T(a,t),\notag\\
 H(n,p)&=\frac{n!\,E_{t+2}^{\,n}}{n^nS}\qquad(n\ge1),
 &H(0,p)&=1.\label{eq:H}
\end{align}
Since $a>6n$, $S\ge1$; no division by zero occurs. From
\eqref{eq:Pierr} and the square-root certificates,
\[
 \abs{S-\sqrt{2\pi n}}\le(2n+1)2^{-t}.
\]
The power factorization and $E_{t+2},e\le3$ give
$\abs{E_{t+2}^n-e^n}\le n3^{n-1}2^{-t}$.
Also $n!\le n^n$, and reciprocals of quantities at least $1$ have
Lipschitz constant at most $1$. Consequently
\begin{equation}\label{eq:Herr}
 \abs{H(n,p)-R_n}
 \le3^n(3n+1)2^{-t}\le2^{-p-2}.
\end{equation}
The last step uses $3^n\le2^{2n}$ and $3n+1\le2^{2n}$ for $n\ge1$,
each a quantifier-free induction. This proves directly that the usual
factorial ratio has a primitive recursive rational name.

\subsection{The relative error and its modulus}

Logarithm addition, $\log e=1$, and
$2\log\sqrt{2\pi n}=\log(2\pi n)$ give
\[
 \log R_n
 =\log(n!)+n-n\log n-\tfrac12\log(2\pi n)=r_n.
\]
Apply the global inequalities \eqref{eq:logglobal} to $R_n$:
\[
 1-1/R_n\le r_n\le R_n-1.
\]
Together with \eqref{eq:strongrobbins}, these imply
\begin{equation}\label{eq:relativeboth}
 \frac1{12n+1}<R_n-1<\frac1{12n-1}.
\end{equation}
For the upper bound, $1-1/R_n<1/(12n)$ and positive-denominator
algebra give $R_n<12n/(12n-1)$. The lower bound follows from
$R_n-1\ge r_n>1/(12n+1)$.
Thus $n\ge2^p$ implies
\[
 \abs{R_n-1}<\frac1{12n-1}\le\frac1n\le2^{-p}.
\]
This completes the proof of Theorem~\ref{thm:main}. The notation
$e^{r_n}$ in the usual logarithmic form of Stirling is the standard
exponential reading of $\log R_n=r_n$; no general exponential function
was needed as an internal premise.

Equations \eqref{eq:Herr} and \eqref{eq:relativeboth} also give the
direct rational arithmetic statement
\begin{equation}\label{eq:Hinternal}
 1+\frac1{12n+1}-2^{-p}\le H(n,p)
 \le1+\frac1{12n-1}+2^{-p}\qquad(n\ge1).
\end{equation}
For $n\ge2^k$, its asymptotic version is
$\abs{H(n,p)-1}\le2^{-k}+2^{-p}$. These are again universal closures
of decidable rational inequalities.

\section{Eliminating the analytic notation inside \texorpdfstring{$\ipra$}{iPRA}}
\label{sec:elimination}

This section records why the proof belongs to the stated arithmetic
theory. Merely having primitive recursive approximants would not establish
provability in $\ipra$; the induction and precision arguments matter.

\subsection{The error of the fixed endpoint \texorpdfstring{$Q(n,p)$}{Q(n,p)}}

Fix $n\ge1,p\ge0$ and abbreviate $b=b(n)$, $t=p+b+4$,
$\delta=2^{-t}$. Every logarithm approximation in \eqref{eq:Q} has
error at most $\delta$. Moreover $2P(t),2\pi>6$, so
\eqref{eq:loglip} and \eqref{eq:Pierr} give
\[
 \tfrac12\abs{\log(2P(t))-\log(2\pi)}\le\delta/6.
\]
The sum of all errors in \eqref{eq:Q} is therefore at most
\begin{align}
 \abs{Q(n,p)-r_n}
 &\le\bigl(n+(n+\tfrac12)+\tfrac12+\tfrac16\bigr)\delta\notag\\
 &= (2n+7/6)\delta
 \le(2n+2)\delta\le2^{-p-3}.\label{eq:Qerr}
\end{align}
The final inequality uses $n+1\le2^b$. This establishes the claimed
name for the logarithmic remainder with considerable slack.

Combining \eqref{eq:finitefinal} with \eqref{eq:Qerr}, for any
$m>n$, gives the rational bounds
\begin{equation}\label{eq:Qfinite}
 g_n+\lmargin_n-\frac1{8m}-2^{-p-3}
 \le Q(n,p)
 \le h_n-\umargin_n+\frac1{6m}+2^{-p-3}.
\end{equation}
Here both endpoints and $Q(n,p)$ are explicit rationals. Choose
$m=\max(n+1,2^{p+8})$. Each total error in \eqref{eq:Qfinite} is
less than $2^{-p}$. We have proved the strengthened arithmetic theorem
\begin{equation}\label{eq:Qstrong}
 \boxed{\quad
 g_n+\lmargin_n-2^{-p}\le Q(n,p)
 \le h_n-\umargin_n+2^{-p}\qquad(n\ge1).
 \quad}
\end{equation}
Dropping the positive margins proves \eqref{eq:internal}.

\subsection{What a finite-precision reading of the proof does}

The preceding named notation is a convenient way of organizing finite
lemmas, not an extra axiom. The following procedure removes it.

\begin{enumerate}
\item Fix the outer parameters $n,p$, the terminal index
$m=\max(n+1,2^{p+8})$, and the tolerances. Do not perform induction
on these choices or on a formula saying an estimate holds at all
precisions.

\item Replace every logarithmic identity by its rational partition proof
from Section~\ref{sec:log}, every series by its displayed finite tail
bound, and every Wallis step by \eqref{eq:finitewallis}. For an auxiliary
tolerance $\eta>0$, a Wallis mesh
$N\ge(D_m+1)/\eta$ suffices. For a logarithm argument $A/B$, a
mesh $N\ge(A+B)^4/\eta$ suffices. Round these bounds up to
integers. For each fixed positive power-series ratio at most $1/3$,
the specified geometric tail gives an equally explicit truncation.

\item Finite sums and products are treated by rational loop invariants.
For example, a sum of $s$ inequalities each with error at most
$\eta/(s+1)$ has total error at most $\eta$. Multiplication by a
rational coefficient $c$ multiplies the error bound by $|c|$.
Division uses its previously proved positive lower bound on the
denominator. All coefficient bounds are rational expressions or finite
sums/products of the parameters already fixed. They are themselves
primitive recursive. Choose the local tolerances by these rules before
starting the corresponding loop.

\item Prove the resulting finite inequalities by quantifier-free
induction, with those tolerances and final loop lengths as parameters.
The finite Wallis induction is explicitly
\eqref{eq:evenerr}--\eqref{eq:odderr}. For logarithmic telescoping use
a common precision: if each instance of \eqref{eq:telescopers} has
error $\eta/(m+1)$, summing at most $m$ instances adds at most
$\eta$. Exact finite-sum cancellation of the chosen rational
approximants is an equation. This proves the finite version of
\eqref{eq:strongtel} without induction on a named-real statement.

\item Allow an extra total error of, say, $2^{-p-4}$ for these
auxiliary replacements. The residual error in \eqref{eq:Qfinite}
is already below $(1/8+1/1536)2^{-p}$, leaving ample room before
$2^{-p}$. Thus one explicit finite precision suffices for
\eqref{eq:Qstrong}. If an intermediate exact rational inequality is
preferred, the elementary lemma below removes its remaining
arbitrarily small auxiliary error.
\end{enumerate}

There is no unbounded search for an adequate stage in this procedure.
Every mesh bound, truncation length, partition refinement, and number of
steps is built from the displayed arithmetic functions and finite sums
or products. A smaller error requirement only changes an explicit natural
parameter.

\begin{lemma}[Removing auxiliary error from a rational comparison]
\label{lem:gap}
Suppose $u,v$ are rational numbers, $D>0$ is a common denominator,
and $B\ge0$ is rational. If finite arithmetic proves
$u\le v+B2^{-k}$ for each natural parameter $k$, then $\ipra$
proves $u\le v$ without a limit principle.
\end{lemma}

\begin{proof}
Choose an integer $M\ge BD$ and put $k=M+1$. The elementary
induction $2^{M+1}>M$ gives $B2^{-k}<1/D$ (the case $B=0$
is immediate). The rational comparison $u\le v$ is decidable. If
it fails, the integer numerator of $u-v$ is positive, so
$u-v\ge1/D$, contradicting this particular finite instance.
This is an intuitionistically valid case split on integers, with an
explicit witness $k$; it is neither Markov's principle nor a use of
real-number trichotomy.
\end{proof}

This lemma justifies, for instance, passing from the finite partition
proofs of logarithm addition to the rational approximation inequalities
used in \eqref{eq:Qfinite}. It never authorizes induction on a
formula quantified over $k$: the finite proof is first carried out with
$k$ as a free parameter, and universal generalization occurs only
afterward.

\subsection{Induction audit}

The relevant invariant types are collected below. In the second column,
all parameters other than the displayed loop index are fixed. Positive
denominators are cleared, and a bounded conjunction is replaced, when
necessary, by a primitive recursive Boolean test.

\begin{longtable}{@{}p{0.28\textwidth}p{0.66\textwidth}@{}}
\toprule
Operation & Quantifier-free invariant used in its proof\\
\midrule
\endhead
Powers, factorials, binomial coefficients & Defining equations and
cross-multiplied finite identities, including Pascal's recurrence.\\[4pt]
Rational partition refinement & Current lower sum does not decrease,
upper sum does not increase; accumulated width is at most the declared
rational budget.\\[4pt]
Geometric and polynomial sums & The finite geometric identity and the
factorization of successive powers, with a finite remainder bound.\\[4pt]
Wallis recurrence & The explicit rational inequalities
\eqref{eq:evenerr}--\eqref{eq:odderr}, with mesh $N$ fixed.\\[4pt]
Logarithmic telescoping & Partial rational sum lies between the two
rational telescopers plus its accumulated precision error.\\[4pt]
Bit length and range reduction & Current largest index within the bounded
scan satisfies the defining power-of-two inequality.\\[4pt]
Square-root search & Current candidate has square at most the fixed
rational input; all previously scanned larger candidates fail. The latter
is a bounded Boolean assertion.\\[4pt]
Precision allocation & A finite sum of local rational errors does not
exceed the assigned total.\\
\bottomrule
\end{longtable}

The only unbounded quantifiers in the final theorem are universal
closures over input and precision variables. The proof does not use
induction on any formula containing those quantifiers. It does not assume
a uniform primitive recursive evaluator for arbitrary primitive recursive
programs: only the finitely many fixed function symbols constructed here
are used. These distinctions explain both the strength and the limits of
the claim that the theorem is proved in $\ipra$.

\section{A concrete path to a complete Lean formalization}
\label{sec:lean}

\subsection{The inspected repository interface}

The following observations refer to the inspected ProveIt revision
\begin{center}\small\repo\end{center}
whose root toolchain and Mathlib dependency are both pinned to version
\texttt{v4.32.0}. The exact source links and inspection scope are in
\texttt{provenance/repository.txt}. The useful part is
\texttt{Logic/PrimitiveRecursiveArithmetic/ConciseIPRA/Lean/ConciseIPRA}.
No result from the Busy Beaver or Diophantine consistency directories is
needed for the numerical proof.

\begin{center}
\begin{tabular}{@{}p{0.40\textwidth}p{0.54\textwidth}@{}}
\toprule
Repository module & Relevant verified declaration or purpose\\
\midrule
\texttt{PR.lean} & Syntax \texttt{PR k}, terms \texttt{PTm n}, axioms
\texttt{PRAAx}, and the rule \texttt{PRA.qfInd}.\\[4pt]
\texttt{Syntax.lean} & First-order formulas \texttt{Fm T n} and their
quantifier-free predicate \texttt{Fm.QF}.\\[4pt]
\texttt{Deduction.lean} & Intuitionistic object-language derivability
\texttt{Prf Ax Gamma phi}.\\[4pt]
\texttt{PRAToScalar.lean} & \texttt{FromPRA.code}, compiling genuine
primitive recursive syntax to scalar iterator terms.\\[4pt]
\texttt{TeqReflect.lean} & Characteristic terms for quantifier-free
formulas and derived quantifier-free induction.\\[4pt]
\texttt{NativeBiInterpretation.lean} & The theorem
\texttt{pra\_iff\_native}, translating an actual iPRA derivation to
the native calculus and back.\\
\bottomrule
\end{tabular}
\end{center}

These interfaces are read from the source, not inferred from filenames
\cite{ProveItPR,ProveItNative}. The archive's Lean scaffold uses them
directly. No Lean or Lake executable was available in the authoring
environment, so the newly supplied Lean source has not been compiled.

\subsection{The mathematical program and the internal theorem differ}

A Lean theorem about \texttt{Rat} or \texttt{Real} is a semantic
mathematical theorem. The requested restriction is an assertion about what
the \emph{object theory} $\ipra$ proves. The final target must therefore
be a derivation with axiom theory \texttt{PRAAx}.

Define a fixed two-input primitive recursive checker that returns $0$
exactly when $n=0$ or the two comparisons in \eqref{eq:internal}
hold. Its object-language open formula is literally
\[
 \ArithCheck(n,p)=0.
\]
The desired closed sentence universally quantifies $p$ and $n$.
The included source defines the formula from a value of type
\texttt{PR 2}; using de Bruijn index $0$ for the innermost variable,
it writes the outer quantifiers in the order $\forall p\,\forall n$.
The order is immaterial to the numerical statement but explicit in the
syntax.

The scaffold separates these two obligations:
\begin{lstlisting}
Representation stirlingApprox
  -- concrete PR syntax and its agreement with the rational algorithm

InternalProof stirlingApprox
  -- the representation plus an actual derivation under PRAAx
\end{lstlisting}
In particular, its \texttt{InternalProof} record has a field of the
object-language derivability type, not merely a field asserting the
standard-model truth of the rational inequality. Constructing this field
is the substantive remaining formalization work. The record is not
instantiated in the delivered scaffold.

Once the derivation exists, universal introduction closes it, and the
existing repository theorem has exactly the required shape:
\begin{lstlisting}
pra_iff_native :
  (Gamma |- [PRAAx] phi) <->
  NPrf (Gamma.map NFm.ofPRA) (NFm.ofPRA phi)
\end{lstlisting}
The ASCII display here is schematic typography. The accompanying Lean
source contains the actual notation and conditional wrapper proofs.

\subsection{Implementation stages and concrete proof obligations}

\paragraph{1. Rational arithmetic in the object language.}
Implement unreduced signed numerator/positive denominator triples as
natural-number codes. Build \texttt{PR} syntax for addition,
multiplication, order, finite sums, powers, and bounded search. Prove the
defining equations and validity invariants under \texttt{PRAAx}.
Use a separate semantic refinement map to Lean's \texttt{Rat} for
readability. Its correctness is useful, but does not replace the internal
arithmetic proofs.

\paragraph{2. A small library of derived quantifier-free proof rules.}
Provide object-language lemmas for transitivity of rational order,
positive cross multiplication, finite sum bounds, product bounds, and
telescoping. Compile bounded universal assertions to Boolean recursors.
Each induction application must display its \texttt{Fm.QF} witness.
For polynomial identities, a normalization procedure should return an
object-language equality derivation, not just a Boolean answer about a
closed numerical example.

\paragraph{3. Finite logarithms.}
Formalize partition refinement and the width
\eqref{eq:logwidth}; prove finite logarithm addition with the sum of
the three enclosure widths. Prove \eqref{eq:monomial} and the
finite geometric remainder \eqref{eq:atanh}. Construct the actual
\texttt{PR} code for \eqref{eq:L}, and derive \eqref{eq:Lerr}
as a family of finite rational comparisons. The supplied source already
defines the corresponding mathematical algorithm over \texttt{Rat}.

\paragraph{4. Finite $\pi$ and Wallis.}
Formalize \eqref{eq:atanaddfinite}, the finite alternating remainder,
and the algorithm \eqref{eq:Pi}. For Wallis, the particularly useful
first machine-checked milestones are \eqref{eq:Jrec},
\eqref{eq:evenerr}, \eqref{eq:odderr}, and
Theorem~\ref{thm:finitewallis}. Their statements involve only
rationals and naturals; their proofs consist of the displayed polynomial
identities, positive denominators, and fixed-precision finite sums.

\paragraph{5. The normalized factorial theorem.}
Prove the coefficient comparisons \eqref{eq:dcompare} and the two
telescopers \eqref{eq:telescopers}. Combine their finite-precision
versions with the factorial identity \eqref{eq:F} and finite Wallis.
Implement the precision schedule of Section~\ref{sec:elimination}.
This produces the open derivation of \eqref{eq:internal}; proving the
stronger \eqref{eq:Qstrong} at the same time records constructive
strictness. Build the \texttt{InternalProof stirlingApprox} value.

\paragraph{6. The direct ratio, if desired.}
Implement the factorial-series and integer-square-root algorithms of
Section~\ref{sec:exp}. Derive \eqref{eq:Hinternal} and its
three-parameter asymptotic version. This optional extension gives a
literal arithmetic approximation to the ratio in the usual formulation,
in addition to the already sufficient logarithmic version.

\paragraph{7. Native transfer and external interpretation.}
Apply the supplied closure and native-transfer wrappers. In a separate
module, interpret the rational names in the ordinary Lean real numbers
and identify them with \texttt{Real.pi}, \texttt{Real.exp 1}, and
the positive square root. This module may use standard real analysis;
the object-language derivation must not depend on it.

Mathlib already contains an ordinary real-analytic Stirling theorem,
including \texttt{Stirling.factorial\_isEquivalent\_stirling}
\cite{MathlibStirling}. It is useful for comparison of statements. Merely
importing that theorem would not demonstrate provability in $\ipra$.

\paragraph{Completion criteria.}
The formalization is complete only when the concrete checker code, its
refinement proof, and its actual \texttt{PRAAx} derivation have been
constructed and checked by Lean. Audit the dependency of the final
derivation on the chosen object theory, and use \texttt{\#print axioms}
on the final Lean declarations as an additional metatheoretic check.
Absence of \texttt{sorry} in a file of definitions and conditional
wrappers is not, by itself, evidence that Stirling has been formalized.

\section{Artifacts, reproducibility, and validation scope}
\label{sec:artifacts}

The archive contains:
\begin{itemize}
\item this LaTeX source and its compiled PDF;
\item \texttt{code/stirling\_rational.py}, an exact-rational reference
implementation of $L$, $P$, $Q$, the finite Wallis sums, and $H$;
\item \texttt{code/run\_checks.py} and
\texttt{validation/checks.json}, recording independently propagated
rational enclosures and finite regression checks;
\item \texttt{lean/ObjectLanguageTarget.lean}, the explicitly
uncompiled concrete algorithms, object-language target, and conditional
native-transfer scaffold;
\item \texttt{provenance/repository.txt}, exact repository references,
and a README with reproduction instructions.
\end{itemize}

The validation run passed \textbf{10,974 exact-rational assertions}.
It includes $Q(n,p)$ for $1\le n\le25$, $0\le p\le12$; direct
ratio cases for $1\le n\le10$, $0\le p\le6$; finite Wallis
recurrences and increment bounds; independently propagated intervals for
logarithms and $\pi$; rational square-root certificates; and
coefficient checks of the proof's polynomial identities. Decimal output
is only for display. These checks helped detect arithmetic or indexing
errors; they do not prove the universal theorem and are not a Lean or
internal-$\ipra$ verification.

To reproduce them from the archive directory, run
\begin{lstlisting}
python3 code/run_checks.py
python3 code/stirling_rational.py 25 12
python3 code/build_pdf.py
\end{lstlisting}
The build helper uses installed TeX files and can generate a private format
and font map when needed. Only the Python standard library is needed for
the numerical checks.
The proof's primitive recursive implementation can use unreduced rational
pairs and bounded integer-square-root search; the reference program uses
Python's exact \texttt{Fraction} and integer arithmetic for convenience.

\begin{thebibliography}{9}
\bibitem{Robbins}
Herbert Robbins, \emph{A remark on Stirling's formula},
The American Mathematical Monthly \textbf{62} (1955), 26--29.
\href{https://heuklyd.github.io/papers/pdf/Robbins-1955.pdf}{Accessible transcription of the paper}.
The classical numerical bounds are attributed to this source; the
finite-rational realization is developed explicitly in this manuscript.

\bibitem{Bruin}
Henk Bruin, \emph{Wallis' Product and Stirling's Formula}, lecture notes.
\url{https://www.mat.univie.ac.at/~bruin/Stirling.pdf}.
Background for the classical relationship between the product and the
normalizing constant.

\bibitem{ProveItPR}
Vladimir Reshetnikov et al., ProveIt, \texttt{ConciseIPRA/PR.lean},
revision \repo.
\href{https://github.com/VladimirReshetnikov/ProveIt/blob/329161f8da0b319d51ba9f6246c4b19d243f7193/Logic/PrimitiveRecursiveArithmetic/ConciseIPRA/Lean/ConciseIPRA/PR.lean}{Pinned source: primitive recursive syntax and iPRA axioms}.

\bibitem{ProveItNative}
ProveIt, \texttt{ConciseIPRA/NativeBiInterpretation.lean}, same revision.
\href{https://github.com/VladimirReshetnikov/ProveIt/blob/329161f8da0b319d51ba9f6246c4b19d243f7193/Logic/PrimitiveRecursiveArithmetic/ConciseIPRA/Lean/ConciseIPRA/NativeBiInterpretation.lean}{Pinned source: native bi-interpretation and transfer theorems}.

\bibitem{MathlibStirling}
The Mathlib contributors, \texttt{Mathlib.Analysis.SpecialFunctions.Stirling},
official documentation, inspected for comparison of the usual real-valued
statement.
\url{https://leanprover-community.github.io/mathlib4_docs/Mathlib/Analysis/SpecialFunctions/Stirling.html}.
The evolving online documentation is not a dependency of the internal
proof; the inspected ProveIt toolchain pin is recorded separately.
\end{thebibliography}
\end{document}


